% Generated draft: requires mathematical and source review.
\documentclass[11pt]{article}
\usepackage[margin=25mm]{geometry}
\usepackage{fontspec}
\setmainfont{Latin Modern Roman}
\usepackage{amsmath,amssymb,mathtools,unicode-math}
\setmathfont{Latin Modern Math}
\usepackage{enumitem,needspace,xurl,hyperref}
\setlength{\parindent}{0pt}
\setlength{\parskip}{5pt}
\setlength{\emergencystretch}{3em}
\newcommand{\N}{\mathbb{N}}
\newcommand{\Z}{\mathbb{Z}}
\newcommand{\Q}{\mathbb{Q}}
\newcommand{\R}{\mathbb{R}}
\newcommand{\C}{\mathbb{C}}
\newcommand{\F}{\mathbb{F}}
\newcommand{\A}{\mathbb{A}}
\newcommand{\PP}{\mathbb{P}}
\newcommand{\E}{\mathbb{E}}
\newcommand{\eps}{\varepsilon}
\newcommand{\dd}{\,\mathrm{d}}
\begin{document}
\section*{UnsolvedMath ID: 800695}
\textbf{Source problem code:} catalogue entry
\subsection{Definitions and mathematical statement}
Let \(F\) be a field. Under the finite-dimensional central interpretation described below, let \(D\) be a division algebra with center \(F\) and \(\dim_F D=n^2<\infty\); its degree is \(\deg D=n\). A field extension \(L/F\) splits \(D\) if
\[
D\otimes_F L\cong M_n(L)
\]
as \(L\)-algebras. Interpret ``split by an abelian group'' as the existence of a finite Galois splitting extension \(L/F\) whose Galois group \(\operatorname{Gal}(L/F)\) is abelian.

With these explicitly stated interpretations, the qualitative question is whether
\[
\begin{gathered}
\text{for every field }F\text{ and every finite-dimensional central division }F\text{-algebra }D,\\
\text{there exists a finite Galois extension }L/F\text{ such that}\\
\operatorname{Gal}(L/F)\text{ is abelian and }D\otimes_F L\cong M_{\deg D}(L).
\end{gathered}
\]
The supplied record also asks for degree bounds over specific fields. It does not identify those fields, specify which degree is to be bounded, or give the parameters on which such a bound may depend. Consequently this quantitative request cannot yet be expressed as a definite mathematical assertion.

The finite-dimensional central setting and the interpretation using an abelian Galois splitting extension are explicit provisional interpretations of the supplied wording, not hypotheses confirmed from the cited source. No characteristic restriction is supplied. The record does not say that \(L\) must embed in \(D\), be a maximal subfield, or satisfy \([L:F]=\deg D\); none of these additional conditions is imposed here.
\subsection{Short English statement}
Under an explicit provisional interpretation, ask whether every finite-dimensional division algebra over its center becomes a matrix algebra after a finite abelian Galois extension. The accompanying request for degree bounds remains unspecified because the record names neither the fields nor the degree to be bounded. This is a statement-only draft requiring human mathematical review.
\subsection{Sources}
\begin{itemize}
\item Local UnsolvedMath registry, permanent ID \(800695\), exact source problem code \texttt{catalogue entry}, title ``Open Problems on Central Simple Algebras --- Problem 5.'' The supplied record attributes the problem to Auel, Brussel, Garibaldi, and Vishne, with proposed year 2010.
\item Supplied reference: Auel, Brussel, Garibaldi, and Vishne, \emph{Open Problems on Central Simple Algebras} (2010), Problem 5, \url{https://arxiv.org/abs/1006.3304}. The background records access on 2026-10-16 and extraction by \texttt{arxiv\_latex}; these are supplied provenance claims, not access or verification performed for this draft.
\item The local record labels the problem \texttt{open} and records an update on 2026-10-16. This is a supplied registry classification, not an independently verified mathematical or literature assessment.
\end{itemize}
\end{document}
